\section{Related Work}

The introduction of Fully Convolutional Neural Networks~\cite{fcn} meant a great advance for the segmentation field. The authors employed an encoder-decoder architecture that only contains convolutional layers. The encoder constitutes the contracting path and the decoder the expansive path, where both are built on repetitive blocks. A landmark framework came with the U-Net architecture~\cite{unet}, where they managed to capture high-level context and local patterns by using skip connections. The underlying idea was to use the same feature maps that are obtained in the contraction to expand the component feature vector. To this aim, the output of each encoder block was stacked to the corresponding input of the decoder block, ensuring that the information that is lost during the downsampling is recovered. 

Distinct variants and applications have come up thanks to the U-Net~\cite{unet}. In V-Net~\cite{vnet}, they presented a deeper model by replacing the stride component by another convolutional layer and a novel loss based on the Dice coefficient to deal with pixel-class imbalances. Attention U-Net~\cite{attentionunet} extended the architecture by an attention grid-based gate for efficient localization. U-Net++~\cite{unet++} redesigned the skip connections with nested convolutional blocks to adapt the intermediate encoder features before stacking them to the decoder block. RefineNet~\cite{refinenet} is a multi-scaled input network that propagates the gradients through the residual connections, enabling the model to obtain high-resolution accurate predictions. Despite the appearance of all these new segmentation techniques, the U-Net~\cite{unet} still demonstrates its capacity to outperform current state-of-the-art techniques in some tasks, such as in satellite image segmentation~\cite{satellite}.

%\cite{hrnet}
%\cite{pathnet}

Other segmentation techniques that also rely on the encoder-decoder framework include a post-processing step to enhance their predicted mapping. DeepLab family~\cite{deeplab} introduced dilated convolutions, capable of enlarging the receptive field of their convolutional without increasing the numbers of weights. In their first work, they utilized fully-connected Conditional Random Fields to fine-tune the localized segmented objects. On the other hand, its later version DeepLab V3~\cite{deeplabv3} achieved encouraging results by combining multiple scales of dilated convolutions through pyramid pooling. DeepLab V3~\cite{deeplabv3} is further strengthened through superpixel merging in~\cite{generic}. Another interesting approach for facilitating the information flow between distinct scales is BPNet~\cite{realtime}. Its architecture is based on a symmetric multi-scale network where the successive layers are continuously merged. Its last stage incorporates an unary attention mechanism for context aggregation. 


%CHALLENGES FOR MY PROBLEM (DATA) AND WHY CURRENT SOLUTIONS DO NOT WORK FOR OUR PROBLEM

%per aixo, proposem un metode de segmentacio per resoldre aquest problemes. la segmentacio es basa en bla bla. Els metodes actuals de segmentacio fallen amb dades noves \cite{}, no son robustos. Per aixo, proposem un metode adaptat, que aprofita el knowledge que tenim de les imatges de pales per tal.